\documentclass[12pt, a4paper]{article}

% ── Paquetes ─────────────────────────────────────────────────────────
\usepackage[utf8]{inputenc}
\usepackage[T1]{fontenc}
% babel-spanish no disponible en esta instalación
\usepackage{amsmath, amssymb, amsthm}
\usepackage{geometry}
\usepackage{graphicx}
\usepackage{xcolor}
\usepackage{mdframed}
\usepackage{enumitem}
\usepackage{booktabs}
\usepackage{hyperref}

% ── Geometría ────────────────────────────────────────────────────────
\geometry{
    top    = 2.5cm,
    bottom = 2.5cm,
    left   = 3.0cm,
    right  = 3.0cm
}

% ── Colores ──────────────────────────────────────────────────────────
\definecolor{azulviz}{RGB}{37, 99, 235}
\definecolor{naranjviz}{RGB}{251, 146, 60}
\definecolor{verdviz}{RGB}{52, 211, 153}
\definecolor{amarillviz}{RGB}{251, 191, 36}
\definecolor{fondocaja}{RGB}{240, 245, 255}
\definecolor{bordecruz}{RGB}{37, 99, 235}
\definecolor{fondoadv}{RGB}{255, 248, 230}
\definecolor{bordeadv}{RGB}{251, 146, 60}
\definecolor{fondoteo}{RGB}{235, 250, 242}
\definecolor{bordeteo}{RGB}{52, 153, 110}

% ── Entornos ─────────────────────────────────────────────────────────
\mdfdefinestyle{vizverde}{
    backgroundcolor = fondoteo,
    linecolor       = bordeteo,
    linewidth       = 1.5pt,
    leftmargin      = 0pt,
    rightmargin     = 0pt,
    innertopmargin  = 8pt,
    innerbottommargin = 8pt,
    innerleftmargin  = 12pt,
    innerrightmargin = 12pt,
    roundcorner     = 4pt
}

\mdfdefinestyle{viznaranja}{
    backgroundcolor = fondoadv,
    linecolor       = bordeadv,
    linewidth       = 1.5pt,
    leftmargin      = 0pt,
    rightmargin     = 0pt,
    innertopmargin  = 8pt,
    innerbottommargin = 8pt,
    innerleftmargin  = 12pt,
    innerrightmargin = 12pt,
    roundcorner     = 4pt
}

\mdfdefinestyle{vizazul}{
    backgroundcolor = fondocaja,
    linecolor       = bordecruz,
    linewidth       = 1.5pt,
    leftmargin      = 0pt,
    rightmargin     = 0pt,
    innertopmargin  = 8pt,
    innerbottommargin = 8pt,
    innerleftmargin  = 12pt,
    innerrightmargin = 12pt,
    roundcorner     = 4pt
}

\newenvironment{viznota}[1]{%
    \begin{mdframed}[style=vizverde]
    \textbf{\color{bordeteo}En GradienViz \textemdash\ #1.}\quad
}{%
    \end{mdframed}
}

\newenvironment{vizadv}{%
    \begin{mdframed}[style=viznaranja]
    \textbf{\color{bordeadv}Atención:}\quad
}{%
    \end{mdframed}
}

\newenvironment{vizdef}{%
    \begin{mdframed}[style=vizazul]
}{%
    \end{mdframed}
}

% ── Comandos matemáticos ─────────────────────────────────────────────
\newcommand{\R}{\mathbb{R}}
\newcommand{\norm}[1]{\left\lVert #1 \right\rVert}
\newcommand{\grad}{\nabla}
\newcommand{\T}{{\top}}
\newcommand{\btheta}{\boldsymbol{\theta}}

% ── Hyperref ─────────────────────────────────────────────────────────
\hypersetup{
    colorlinks = true,
    linkcolor  = azulviz,
    urlcolor   = azulviz,
    citecolor  = azulviz
}

% ── Metadatos ────────────────────────────────────────────────────────
\title{
    {\Large\bfseries Más allá de la regresión lineal}\\[0.4em]
    {\large Lo que cambia cuando el modelo cambia}\\[0.8em]
    {\normalsize Notas de lectura para \textit{GradienViz}}
}
\author{
    Departamento Académico de Sistemas Computacionales\\
    Universidad Autónoma de Baja California Sur
}
\date{2026}

% ─────────────────────────────────────────────────────────────────────
\begin{document}

\maketitle
\tableofcontents
\vspace{1em}

\begin{vizadv}
Este documento complementa los otros tres de la colección GradienViz.
Se asume familiaridad con la regresión lineal simple y con la idea
general del descenso por gradiente.
\end{vizadv}

% =============================================================================
\section{El límite del ejemplo en clase}
% =============================================================================

La regresión lineal simple\textemdash ajustar una recta $y = mx + b$
a un conjunto de puntos\textemdash es el punto de partida ideal para
introducir la optimización. El modelo tiene exactamente dos parámetros,
la función de costo es un paraboloide perfectamente convexo, y existe
una fórmula analítica que lo resuelve en un paso.

Pero los datos reales rara vez son lineales.

\begin{itemize}
    \item La concentración de una bacteria en un cultivo crece de forma
    \textbf{exponencial} con el tiempo: $N(t) = N_0\,e^{rt}$.

    \item La magnitud de un terremoto y la frecuencia con que ocurre
    siguen una \textbf{ley de potencias}: $\log F = a - b\,M$
    (ley de Gutenberg--Richter).

    \item El precio de una vivienda depende simultáneamente de su
    superficie, número de habitaciones, distancia al centro y decenas
    de otras variables: \textbf{regresión múltiple}.

    \item La probabilidad de que un correo sea spam dado su contenido
    sigue un \textbf{modelo logístico}: $P = 1/(1 + e^{-\btheta^\T \mathbf{x}})$.
\end{itemize}

La pregunta que organiza este documento es: \textbf{¿qué cambia en el
método cuando el modelo cambia?} La respuesta, según el nivel de análisis,
va de ``todo'' a ``nada''.

% =============================================================================
\section{Modelos que se linealizan}
% =============================================================================

Una familia amplia de modelos no lineales admite una
\textbf{linealización}: una transformación algebraica que los convierte
en un problema de regresión lineal sobre variables nuevas.
En esos casos, toda la maquinaria del álgebra lineal\textemdash
ecuaciones normales, número de condición, descenso por gradiente\textemdash
se reutiliza sin modificación.

\subsection{Modelo exponencial}

Supón que los datos sugieren $y = a\,e^{bx}$. Aplicando logaritmo natural:
\[
    \ln y = \ln a + b\,x.
\]
Definiendo $y' = \ln y$, $\alpha = \ln a$ y $\beta = b$, el modelo se
convierte en:
\[
    y' = \alpha + \beta x,
\]
que es regresión lineal sobre las variables transformadas $(x,\, \ln y)$.
La matriz de diseño y las ecuaciones normales son idénticas al caso base:
\[
    \mathbf{A} = \begin{pmatrix} x_1 & 1 \\ \vdots & \vdots \\ x_n & 1 \end{pmatrix},
    \qquad
    \mathbf{A}^\T\mathbf{A}\,\btheta^* = \mathbf{A}^\T\mathbf{b}',
    \quad \btheta^* = (\beta,\, \alpha)^\T.
\]
Una vez obtenidos $\alpha$ y $\beta$, se recuperan los parámetros
originales: $b = \beta$, $a = e^{\alpha}$.

\subsection{Modelo potencial}

El modelo $y = a\,x^b$ (frecuente en alometría, sismología, economía)
se linealiza aplicando logaritmo a ambos lados:
\[
    \ln y = \ln a + b\,\ln x.
\]
Ahora tanto $x$ como $y$ aparecen en escala logarítmica. La variable
predictora pasa a ser $x' = \ln x$ y la respuesta $y' = \ln y$:
misma matriz $\mathbf{A}$, misma ecuación normal, dos parámetros.

\subsection{Regresión polinomial}

El modelo $y = \beta_0 + \beta_1 x + \beta_2 x^2 + \cdots + \beta_k x^k$
parece no lineal, pero \emph{es lineal en los parámetros} $\beta_j$.
Basta agregar columnas a $\mathbf{A}$:
\[
    \mathbf{A} = \begin{pmatrix}
        1 & x_1 & x_1^2 & \cdots & x_1^k \\
        1 & x_2 & x_2^2 & \cdots & x_2^k \\
        \vdots & & & & \vdots \\
        1 & x_n & x_n^2 & \cdots & x_n^k
    \end{pmatrix} \in \R^{n \times (k+1)}.
\]
La ecuación normal $\mathbf{A}^\T\mathbf{A}\,\btheta^* = \mathbf{A}^\T\mathbf{b}$
sigue siendo la misma expresión; ahora $\btheta^* \in \R^{k+1}$.

\subsection{Regresión múltiple}

Con $p$ variables predictoras $x_1, \ldots, x_p$, el modelo es:
\[
    y = \beta_0 + \beta_1 x_1 + \beta_2 x_2 + \cdots + \beta_p x_p.
\]
La matriz de diseño crece en columnas:
\[
    \mathbf{A} = \begin{pmatrix}
        1 & x_{11} & x_{12} & \cdots & x_{1p} \\
        1 & x_{21} & x_{22} & \cdots & x_{2p} \\
        \vdots & & & & \vdots \\
        1 & x_{n1} & x_{n2} & \cdots & x_{np}
    \end{pmatrix} \in \R^{n \times (p+1)},
\]
y la ecuación normal permanece inalterada: $\mathbf{A}^\T\mathbf{A}\,\btheta^* = \mathbf{A}^\T\mathbf{b}$,
con $\btheta^* \in \R^{p+1}$.

\begin{vizdef}
\textbf{Patrón unificador.} En todos los modelos linealizables, el método
se reduce a elegir las columnas adecuadas para $\mathbf{A}$ (variables originales,
transformadas, potencias, productos cruzados) y resolver \emph{siempre la
misma ecuación normal}. El álgebra lineal absorbe la generalización; el
cálculo tendría que rehacer todo el análisis para cada modelo nuevo.
\end{vizdef}

\subsection{Lo que sí cambia: la dimensión y el condicionamiento}

Al pasar de 2 parámetros a $p+1$, el espacio de parámetros pasa de
$\R^2$ a $\R^{p+1}$: ya no es visualizable directamente. Lo que
GradienViz muestra en el Panel~2 como un plano con un punto $\oplus$
se convierte en un hiperplano en alta dimensión con un punto mínimo.

La geometría del bowl\textemdash convexa, con un único mínimo global\textemdash
se conserva, pero el número de condición $\kappa(\mathbf{A}^\T\mathbf{A})$
puede crecer drásticamente cuando las columnas de $\mathbf{A}$ están
correlacionadas (multicolinealidad). El Panel~1 de GradienViz ilustra
con dos variables la misma degradación que ocurre en alta dimensión.

\begin{viznota}{Panel~2}
El heatmap y las trayectorias que ves son una ventana al caso $p=1$.
Para regresión múltiple con $p > 2$, la visualización directa es imposible,
pero el bowl de $J$ sigue siendo cuadrático y convexo. Lo que cambia
es solo el número de columnas de $\mathbf{A}$, no la ecuación que lo define.
\end{viznota}

% =============================================================================
\section{Modelos que no se linealizan}
% =============================================================================

No todos los modelos admiten linealización. Los dos casos más importantes
en aprendizaje automático son:

\textbf{Regresión logística.} La probabilidad de pertenecer a una clase
se modela como:
\[
    P(y=1 \mid \mathbf{x}) = \sigma(\btheta^\T \mathbf{x})
    = \frac{1}{1 + e^{-\btheta^\T \mathbf{x}}}.
\]
La función de costo (entropía cruzada) no es cuadrática en $\btheta$.
No existe fórmula cerrada análoga a las ecuaciones normales. El descenso
por gradiente deja de ser pedagógicamente motivado y se convierte en el
\emph{único} método práctico disponible.

\textbf{Redes neuronales.} Componen múltiples transformaciones no lineales.
La función de costo tiene mínimos locales, puntos de silla y regiones
planas extendidas. La ecuación normal no existe. El descenso por gradiente\textemdash
con el gradiente calculado mediante retropropagación\textemdash es el
algoritmo universal de entrenamiento.

\begin{vizadv}
La frontera entre modelos linealizables y no linealizables no es estética:
define si existe o no una solución analítica. En el lado linealizable,
el descenso por gradiente es una elección; en el lado no linealizable,
es una necesidad.
\end{vizadv}

% =============================================================================
\section{Lo que no cambia: la iteración}
% =============================================================================

A pesar de toda la variedad de modelos, funciones de costo y espacios
de parámetros, la regla de actualización del descenso por gradiente es
siempre la misma:

\begin{vizdef}
\textbf{Regla universal.} Sea $\btheta \in \R^p$ el vector de parámetros
del modelo y $J(\btheta)$ cualquier función de costo diferenciable.
La iteración de descenso por gradiente es:
\begin{equation}\label{eq:universal}
    \btheta^{(t+1)} = \btheta^{(t)} - \eta\,\grad J\!\left(\btheta^{(t)}\right),
\end{equation}
independientemente de si el modelo es lineal, exponencial, polinomial,
logístico o una red neuronal profunda.
\end{vizdef}

Lo que varía de un modelo a otro es únicamente la expresión de
$\grad J(\btheta)$: cómo se calcula ese vector en cada caso.
Para modelos lineales es una multiplicación de matrices; para redes
neuronales, la retropropagación. Pero la forma de~\eqref{eq:universal}
no se toca.

\begin{center}
\small
\begin{tabular}{llll}
    \toprule
    \textbf{Modelo} & \textbf{Parámetros} & \textbf{¿Linealizable?} & \textbf{$\grad J$} \\
    \midrule
    Lineal $y = mx + b$          & 2        & —          & Fórmula directa \\
    Exponencial $y = ae^{bx}$    & 2        & Sí ($\ln$) & Fórmula directa \\
    Potencial $y = ax^b$         & 2        & Sí ($\ln$) & Fórmula directa \\
    Polinomial grado $k$         & $k+1$    & Sí         & Fórmula directa \\
    Múltiple $p$ predictores     & $p+1$    & Sí         & Multiplicación matricial \\
    Logístico                    & $p+1$    & No         & Fórmula directa \\
    Red neuronal $2\to 4\to 1$   & 17       & No         & Retropropagación \\
    Red neuronal profunda        & $10^{9}$ & No         & Retropropagación \\
    \bottomrule
\end{tabular}
\end{center}

\bigskip

GradienViz muestra la fila más sencilla de esa tabla. Lo que el
visualizador hace visible\textemdash trayectorias que siguen el gradiente
negativo hacia un mínimo\textemdash es el mismo proceso que ocurre,
invisible y en miles de millones de dimensiones, cada vez que se entrena
un modelo de lenguaje, un sistema de reconocimiento de imágenes o un
motor de recomendación.

La regresión lineal no es el destino. Es la única situación donde el
destino se puede ver.

\end{document}
